% =====================================================================
% CHAPTER 5: METHODOLOGY
% =====================================================================
\chapter{Methodology}
\label{ch:methodology}

This chapter documents the complete implemented pipeline of TrackVision, from video acquisition to visualization. Every subsection corresponds to actual code in the repository.

\section{System Overview}
\label{sec:sys-overview}

TrackVision follows a \textbf{tracking-by-detection} architecture with four Web Workers operating in parallel, coordinated by a main-thread vision engine loop (requestAnimationFrame). The high-level data flow is:

\begin{center}
\texttt{Video} $\to$ \texttt{OffscreenCanvas} $\to$ \texttt{Detection Workers} $\to$ \texttt{ReID Worker} $\to$ \texttt{Tracker Worker} $\to$ \texttt{State Store} $\to$ \texttt{UI}
\end{center}

Figure~\ref{fig:architecture} illustrates the worker architecture.

\section{Browser Video Acquisition}
\label{sec:video-acq}

Video input is obtained via the standard \texttt{MediaDevices API}:

\begin{lstlisting}[style=typescript, caption={Camera request in CameraHUD.tsx}]
navigator.mediaDevices.getUserMedia({
  video: { 
    facingMode: 'environment', 
    width: { ideal: 1280 }, 
    height: { ideal: 720 } 
  },
  audio: false
}).then(stream => {
  videoRef.current.srcObject = stream;
  videoRef.current.onloadedmetadata = () => videoRef.current?.play();
});
\end{lstlisting}

The rear-facing camera (\texttt{facingMode: 'environment'}) is requested by default for mobile scenarios. On desktop, the user's default camera is used. Camera permission is required; errors surface as UI banners with HTTPS guidance for insecure origins.

\section{Frame Capture and Letterboxing}
\label{sec:frame-capture}

Frames are captured using \texttt{OffscreenCanvas} for zero-main-thread overhead:

\begin{lstlisting}[style=typescript, caption={Frame capture in useOffscreenCanvas.ts}]
captureFrame(targetW: number, targetH: number): CapturedFrame | null {
  const canvas = offscreenCanvasRef.current;
  const video = videoRef.current;
  if (!canvas || !video || video.readyState < 2) return null;

  const scale = Math.min(targetW / video.videoWidth, targetH / video.videoHeight);
  const newW = Math.round(video.videoWidth * scale);
  const newH = Math.round(video.videoHeight * scale);
  const offsetX = (targetW - newW) / 2;
  const offsetY = (targetH - newH) / 2;

  const ctx = canvas.getContext('2d');
  ctx.clearRect(0, 0, targetW, targetH);
  ctx.drawImage(video, offsetX, offsetY, newW, newH);

  const imageData = ctx.getImageData(0, 0, targetW, targetH);
  const buffer = imageData.data.buffer.slice(0); // copy for transfer

  return { buffer, data: imageData.data, width: targetW, height: targetH,
           originalWidth: video.videoWidth, originalHeight: video.videoHeight,
           scale, offsetX, offsetY };
}
\end{lstlisting}

Key properties:
\begin{itemize}
    \item \textbf{Letterboxing:} Aspect ratio preserved; padding equally distributed.
    \item \textbf{Metadata:} \texttt{scale}, \texttt{offsetX}, \texttt{offsetY} enable mapping detection coordinates back to original video space.
    \item \textbf{Transfer:} \texttt{ArrayBuffer} is transferred (detached) to workers via \texttt{postMessage}, avoiding copies.
\end{itemize}

\section{Preprocessing for Inference}
\label{sec:preprocessing}

Each worker receives the transferred buffer and performs planar conversion:

\begin{lstlisting}[style=typescript, caption={Planar conversion in yoloDetectionWorker.ts}]
const u8 = new Uint8Array(buffer);
const inv255 = 1 / 255;
for (let i = 0, p = 0; i < numPixels; i++, p += 4) {
  planarData[rOffset + i] = u8[p] * inv255;
  planarData[gOffset + i] = u8[p + 1] * inv255;
  planarData[bOffset + i] = u8[p + 2] * inv255;
}
\end{lstlisting}

For ReID (OSNet), crops are extracted from the letterboxed frame, resized to 256x128, and normalized with ImageNet statistics (mean=[0.485, 0.456, 0.406], std=[0.229, 0.224, 0.225]).

\section{Object Detection}
\label{sec:detection}

Two detection paths are implemented:

\subsection{Fast Mode: YOLOv8n}
\label{sec:yolo-fast}
\begin{itemize}
    \item Model: \texttt{yolov8n.onnx} (12.8 MB, COCO-80)
    \item Worker: \texttt{yoloDetectionWorker.ts}
    \item Input: \texttt{float32[1, 3, 640, 640]}
    \item Output: Transposed \texttt{[1, 84, 8400]} (84 = 4 bbox + 80 classes) or non-transposed
    \item Postprocessing: Sigmoid calibration, class-aware NMS (IoU 0.45), coordinate mapping via \texttt{scale/offsetX/offsetY}
\end{itemize}

\subsection{Open-Vocabulary Mode: YOLO-World + CLIP}
\label{sec:yolo-open}
\begin{itemize}
    \item Model: \texttt{yoloworld.onnx} (~418 MB) + \texttt{clip\_text\_encoder.onnx} (~254 MB)
    \item Worker: \texttt{yoloWorldWorker.ts}
    \item CLIP Tokenizer: \texttt{clipTokenizer.ts} (BPE, vocab/merges from \texttt{public/models/clip-tokenizer/})
    \item Text Encoding: \texttt{``a photo of a \{concept\}''} $\to$ tokenize (int64[77]) $\to$ CLIP encoder (WASM) $\to$ L2-normalize $\to$ \texttt{text\_features [1, C, 512]}
    \item YOLO-World Input: \texttt{images [1, 3, 640, 640]} + \texttt{text\_features [1, C, 512]}
    \item Output: Post-sigmoid \texttt{scores [1, N, C]} + XYXY \texttt{boxes [1, N, 4]}
    \item Postprocessing: \texttt{postprocessWorld} decodes scores/boxes in letterbox space
    \item Lazy Loading: CLIP loads on first concept set; embeddings computed once per concept set
\end{itemize}

\section{Detection Filtering and Class-Aware NMS}
\label{sec:nms}

After detection, \textbf{class-aware NMS} runs on the main thread (\texttt{applyNMSWithClass} in \texttt{workerUtils.ts}):

\begin{enumerate}
    \item Group detections by \texttt{classId}.
    \item Within each class, sort by confidence descending.
    \item Iterate: keep highest, suppress others with IoU $> 0.45$.
    \item Merge results across classes.
\end{enumerate}

This prevents cross-class suppression (e.g., a high-confidence ``person'' does not suppress a nearby ``laptop'').

\section{ReID Feature Extraction}
\label{sec:reid}

ReID is \textbf{conditional} -- only active when a valid OSNet model is loaded (\texttt{reidReady \&\& !reidIsFallback}).

\begin{enumerate}
    \item \textbf{Crop Mapping:} Video-space boxes $\to$ letterbox coordinates using \texttt{scale/offsetX/offsetY}.
    \item \textbf{Batch-1 Inference:} Each crop processed sequentially (OSNet export rejects stacked batches).
    \item \textbf{Preprocessing:} Resize 256x128, ImageNet mean/std, RGB $\to$ CHW.
    \item \textbf{Inference:} \texttt{extractBatch} in \texttt{reidWorker.ts} returns 512-dim L2-normalized embedding per crop.
    \item \textbf{Attachment:} Non-zero embeddings attached to detections if length matches.
\end{enumerate}

If OSNet is unavailable, the tracker falls back to IoU + center distance + class consistency only.

\section{Track Prediction}
\label{sec:track-pred}

Before association, every active track predicts its next state via Kalman filter:

\begin{lstlisting}[style=typescript, caption={Track prediction in tracker.ts}]
predict(): BBox {
  const predicted = this.kalman.predict();
  this.bbox = predicted;
  return predicted;
}
\end{lstlisting}

The Kalman filter (\texttt{kalman.ts}) maintains a 6x6 covariance matrix with constant acceleration process noise.

\section{Data Association}
\label{sec:association}

The core of ByteTrack++ is the \textbf{four-stage association} (\texttt{ByteTracker.update} in \texttt{tracker.ts}):

\subsection{Stage 1: High-Confidence Association}
Match detections with $\text{score} \ge \text{trackThresh}$ (capped at 0.4) against all existing tracks using Hungarian algorithm on combined cost matrix.

\subsection{Stage 2: Low-Confidence Association}
Match detections with $0.1 < \text{score} < \text{trackThresh}$ against only the \textit{unmatched} tracks from Stage 1.

\subsection{Stage 3: New Track Creation}
Unmatched high-confidence detections spawn new tracks (TENTATIVE state, new Kalman filter).

\subsection{Stage 4: Missed Track Handling}
Unmatched tracks increment \texttt{timeSinceUpdate}; if $> 5$ frames $\to$ LOST; if $> 30$ frames (\texttt{maxTimeLost}) $\to$ REMOVED.

\subsection{Cost Matrix Composition}
\label{sec:cost-matrix}
For each track-detection pair, the combined cost is:

\[
\text{cost} = 
\begin{cases}
0.6 \cdot (1 - \text{IoU}) + 0.3 \cdot (1 - \text{cosine\_sim}) + 0.1 \cdot \text{class\_cost} + 0.1 \cdot \text{center\_cost} & \text{if embeddings exist} \\
0.8 \cdot (1 - \text{IoU}) + 0.15 \cdot \text{center\_cost} + 0.05 \cdot \text{class\_cost} & \text{otherwise}
\end{cases}
\]

Where:
\begin{itemize}
    \item $\text{class\_cost} = 0$ if same class, $0.4$ otherwise
    \item $\text{center\_cost} = \min(1, \frac{\text{center\_dist}}{2 \cdot \max(\text{dims})})$
    \item Gating: costs $> 1 - \text{matchThresh}$ (default 0.3) set to $\infty$ before Hungarian
\end{itemize}

\section{Track Lifecycle Management}
\label{sec:lifecycle}

Tracks transition through states:

\begin{center}
\texttt{TENTATIVE} $\xrightarrow{\text{hits} \ge 3 \text{ or } \text{confirmedHits} \ge 2}$ \texttt{CONFIRMED} $\xrightarrow{\text{timeSinceUpdate} > 5}$ \texttt{LOST} $\xrightarrow{\text{timeSinceUpdate} > 30}$ \texttt{REMOVED}
\end{center}

\begin{itemize}
    \item \texttt{confirmedHits}: detections with $\text{score} \ge 0.7$
    \item Recovery: LOST $\to$ TRACKED if re-associated within \texttt{maxTimeLost}
    \item Label stabilization: Majority vote over 10 class labels
    \item Score smoothing: EMA $\alpha=0.3$ ($\text{score} \leftarrow 0.7 \cdot \text{old} + 0.3 \cdot \text{new}$)
\end{itemize}

\section{Track Merging (Split Identity Recovery)}
\label{sec:merging}

If two confirmed tracks of the \textit{same class} have $\text{IoU} > 0.5$ and appearance similarity $> 0.7$, the older track absorbs the younger (EMA embedding fusion, younger marked REMOVED).

\section{EMA Smoothing and Interpolation}
\label{sec:smoothing}

\begin{itemize}
    \item \textbf{BBox EMA:} $\alpha=0.7$ on \texttt{bbox} coordinates.
    \item \textbf{Center EMA:} $\alpha=0.8$ on center points.
    \item \textbf{Occlusion Interpolation:} Up to 5 frames of Kalman-predicted positions blended with smoothed bbox when \texttt{timeSinceUpdate} $> 0$.
\end{itemize}

\section{Output Rendering}
\label{sec:rendering}

The main thread renders tracking overlays on a canvas (\texttt{drawTrackingHUD} in \texttt{lib/draw.ts}):

\begin{itemize}
    \item Bounding boxes with track IDs, class labels, confidence.
    \item Ghost mode: Fading trajectory trails (15 frames, opacity decay).
    \item Follow mode: Camera locked to selected track, dimmed background.
    \item Scene Map: 640x480 top-down density map with live track dots.
    \item Timeline: Gantt-chart horizontal bars per track ID (5-min window).
\end{itemize}

\section{Telemetry Collection}
\label{sec:telemetry}

Per-frame latency breakdown collected in \texttt{useVisionEngine.ts}:

\begin{itemize}
    \item \texttt{preprocessMs}, \texttt{inferenceMs}, \texttt{postprocessMs}, \texttt{trackingMs}, \texttt{renderMs}, \texttt{frameMs}
    \item \texttt{executionProvider} (WebNN/WebGPU/WebGL/WASM)
    \item \texttt{fps} (rolling 500ms window)
    \item History: 60-sample rolling buffers for latency/FPS charts
\end{itemize}

\section{Five-Minute Rolling History}
\label{sec:history}

The Zustand store (\texttt{store.ts}) maintains a sliding window:

\begin{lstlisting}[style=typescript, caption={History pruning in store.ts}]
const cutoffTime = frame.timestamp - state.maxHistoryMs; // 5 * 60 * 1000
while (newFrames.length > 0 && newFrames[0].timestamp < cutoffTime) {
  newFrames.shift();
}
\end{lstlisting}

Track metadata (trajectories, distance, speed, bearing) is updated per frame and pruned to 300 history points per track.

\section{Replay Mode}
\label{sec:replay}

When \texttt{isReplay=true}, the live loop pauses. \texttt{currentTime} scrubs through stored \texttt{frames[]}, rendering historical frames with full track state reconstruction.